{\looseness=-1
\chapter{实验体会}\label{part:experience}

历时数周完成《数据分析与可视化实践》“知微”可视分析系统，从最初对着二十多万条异构 JSONL 原始日志发愁，到如今能够一键启动标准库本地服务、在五球时空入口和七大核心图版之间流畅切换并完成全链路回归测试，这段经历让我们深切体会到数据工程的严谨性与可视设计的克制美感。

\textbf{首先，数据口径的真实与干净，是一切漂亮图表的立身之本。}在项目初期，我们曾直观地以为把几个实验的 CSV 拼在一起就能画出漂亮的对比图。但真正动手清洗时才发现，代码生成里有 465 条因评测环境依赖无法判定的记录，文档位置实验和文档提示实验的采集日期存在跨度，代码修复中六种策略还共享着同一份 Round 0 初稿。若当时把缺失填零，或把共享初稿重复计入物理分母，后续所有置信区间与修复率都会失真。我们在 \texttt{classify\_and\_audit.py} 和 \texttt{quality\_report.json} 上花掉了大量反复核对的时间，甚至为了区分历史 Mini 批次和补充批次特意维护了两套独立口径。如今回头看，正是这份对每一个分母、每一处 \texttt{null} 语义的较真，才让团队在撰写结论时敢于直面“9/12 组限长反弹”或“GSM8K 中文微降”这些不完美但完全真实可验的实证结论。

\textbf{其次，好的可视交互系统应当服务于“发现机制”，而不是堆砌花哨控件。}从早期的 \texttt{explore.html} 到后来重构的“知微”五屏架构，我们推翻过几版设计。最初我们直接调用通用图表库铺满大屏，但很快就发现常规模板很难把“策略矩阵—指令文本—热力矩阵”或者“六条修复轨迹 + 真实代码行级 Diff 抽屉”有机串联在一起。下定决心在 \texttt{zhiwei.js} 中用原生 SVG 手写坐标系、误差线与动态高亮逻辑后，虽然代码量增加到了近两千行，还要处理宽屏吸附、窄屏堆叠和键盘无障碍焦点等繁琐细节，但当鼠标悬停在“中文作答”上、左侧真实指令与右侧 12 个模型得分格同时点亮的那一刻，整个系统真正“活”了起来——图不再是文章的附庸，而成为了读者亲手核验结论的交互实验室。

\textbf{再者，四人并肩攻关的默契，是我们走通工程闭环的最大底气。}面对数据清洗、聚类 Bootstrap 复算、前端原生渲染和浏览器自动化测试四条并行的战线，大家在约定好 21 列 Schema 和四份科研 JSON 契约后便分头推进。每当遇到难以排查的问题——例如 \texttt{build\_web\_app.py} 校验哈希不匹配、或是无头浏览器在 $390\times844$ 手机视口下出现 2 像素横向溢出时，四个人总会在深夜把测试截图和控制台报错发到群里一起定位：有人核对底层样本分母，有人调整 SVG \texttt{viewBox} 边距，有人补写防泄漏单元测试，还有人逐页检查文档链接和表格排版。原本分头推进的脚本与界面，正是在这一次次联调验证中咬合成了完整的系统。

\par}
\textbf{最后，感谢指导老师的启发，也感谢这段时间里始终保持耐心的每一位队友。}这门课程让我们明白，数据分析与可视化要求研究者既守住数据边界与伦理底线，又把复杂规律转化为清晰、诚实、可交互的视觉语言。这份对数据边界与工程质量的敬畏，也将伴随我们走向今后的研发实践。
